% =====================================================================
%  Thème 5 — Angle et identités remarquables — DIFFICILE — Exercices
% =====================================================================
\documentclass[11pt,a4paper]{article}
\usepackage[utf8]{inputenc}
\usepackage[T1]{fontenc}
\usepackage{lmodern}
\usepackage[french]{babel}
\usepackage{amsmath,amssymb,mathtools}
\usepackage{xcolor}
\usepackage{enumitem}
\usepackage{fancyhdr}
\usepackage[a4paper,margin=2.1cm,headheight=26pt,headsep=10pt]{geometry}

\definecolor{primary}{RGB}{222,70,80}
\definecolor{primarydark}{RGB}{150,30,42}

\pagestyle{fancy}\fancyhf{}
\fancyhead[L]{\small\textcolor{primarydark}{\textbf{Fiche 2} — Calcul vectoriel}}
\fancyhead[R]{\small\textcolor{primarydark}{\textsc{Thème 5 — Difficile — Exercices}}}
\fancyfoot[C]{\small\thepage}
\renewcommand{\headrulewidth}{0.8pt}

\newcommand{\vc}[1]{\overrightarrow{#1}}
\providecommand{\norm}[1]{\left\lVert #1\right\rVert}
\newcommand{\exo}[1]{\medskip\noindent{\color{primary}\bfseries\large Exercice #1.}\par\nopagebreak\smallskip}

\begin{document}
\begin{center}
{\LARGE\bfseries\color{primarydark}Angle et identités remarquables — Niveau Difficile}\\[2pt]
{\color{primary}\rule{0.55\linewidth}{1.1pt}}
\end{center}
\vspace{1mm}
{\small\itshape Objectif : remonter d'un angle imposé vers une inconnue, mener une étude paramétrée
complète et conduire une synthèse « vrai / faux » de type concours.}
\vspace{2mm}

\exo{1} \textbf{Retrouver une inconnue à partir d'un angle.}
Soit $\vec u\,(1;x)$ et $\vec v\,(\sqrt3;1)$, où $x$ est un réel \emph{strictement positif}.
On veut que l'angle entre $\vec u$ et $\vec v$ vaille $30^\circ$ (soit $\tfrac{\pi}{6}$).
\begin{enumerate}[label=\textbf{\alph*)},leftmargin=2em,itemsep=3pt]
  \item Exprimer $\vec u\cdot\vec v$, $\norm{\vec u}$ et $\norm{\vec v}$.
  \item Écrire l'équation $\cos30^\circ=\dfrac{\vec u\cdot\vec v}{\norm{\vec u}\,\norm{\vec v}}$
        et l'amener à la forme $\sqrt3+x=\sqrt3\,\sqrt{1+x^2}$.
  \item Résoudre (en élevant au carré) et déterminer $x$.
\end{enumerate}

\exo{2} \textbf{Étude paramétrée complète.}
On donne $\vec u\,(a;a)$ et $\vec v\,\bigl(a;\tfrac{a}{3}\bigr)$, avec $a>0$.
\begin{enumerate}[label=\textbf{\alph*)},leftmargin=2em,itemsep=3pt]
  \item Calculer $\norm{\vec u}$ et $\norm{\vec v}$ en fonction de $a$.
  \item Calculer $\vec u\cdot\vec v$ en fonction de $a$.
  \item En déduire $\cos(\vec u,\vec v)$ et vérifier qu'il \emph{ne dépend pas} de $a$.
  \item Donner une valeur approchée de l'angle (au dixième de degré).
\end{enumerate}

\exo{3} \textbf{Identités remarquables et loi du parallélogramme.}
On donne $\vec u\,(3;1)$ et $\vec v\,(1;2)$.
\begin{enumerate}[label=\textbf{\alph*)},leftmargin=2em,itemsep=3pt]
  \item Calculer $\norm{\vec u}^2$, $\norm{\vec v}^2$ et $\vec u\cdot\vec v$.
  \item À l'aide des identités remarquables, calculer $\norm{\vec u+\vec v}^2$ et $\norm{\vec u-\vec v}^2$.
  \item Vérifier l'\emph{égalité du parallélogramme} :
        $\norm{\vec u+\vec v}^2+\norm{\vec u-\vec v}^2=2\bigl(\norm{\vec u}^2+\norm{\vec v}^2\bigr)$.
\end{enumerate}

\exo{4} \textbf{Synthèse « vrai / faux » (type concours).}
On donne $\vec u\,(2;\sqrt5)$ et $\vec v\,(-4;0)$. Répondre par \textbf{Vrai} ou \textbf{Faux} en
justifiant.
\begin{enumerate}[label=\textbf{\Alph*.},leftmargin=2.2em,itemsep=3pt]
  \item $\norm{\vec u}=3$.
  \item Le produit scalaire $\vec u\cdot\vec v$ est un nombre négatif.
  \item $\norm{\vec u-\vec v}^2=41$.
  \item L'angle entre $\vec u$ et $\vec v$ est aigu.
\end{enumerate}

\end{document}
